\chapter{General Introduction}

\section{Context}
Healthcare systems are increasingly supported by digital technologies such as connected medical devices, electronic health records, telemedicine platforms, mobile health applications, and artificial intelligence-based analytics. This evolution has contributed to the emergence of eHealth, a broad domain that aims to improve healthcare delivery, monitoring, prevention, diagnosis, and patient engagement through information and communication technologies.

At the same time, healthcare data are becoming more heterogeneous and continuous. A patient may generate information through clinical consultations, laboratory tests, medical imaging, wearable devices, mobile applications, and hospital information systems. The challenge is no longer only to collect data, but also to transform these data into meaningful, reliable, and actionable knowledge for healthcare professionals and patients.

Digital Twin technology has recently gained attention as a promising paradigm for representing physical entities and processes through dynamic virtual counterparts. In healthcare, a Digital Twin may support the representation of a patient, a medical device, a hospital service, a clinical workflow, or even a population-level process. The objective is to connect real-world data with virtual models in order to monitor current states, simulate possible futures, and support decision-making.

\section{Problem Statement}
Despite the progress of eHealth systems, several limitations remain. First, medical data are often distributed across multiple systems and formats, which makes integration and interoperability difficult. Second, many digital health applications provide static analysis based on isolated measurements rather than dynamic representations that evolve with the real-world situation. Third, decision support systems frequently rely on general population models and may not sufficiently account for individual variability, contextual information, and temporal evolution.

These limitations raise the following central problem:

\begin{quote}
How can Digital Twin technology be used as a conceptual and architectural foundation for building eHealth systems that are more dynamic, interoperable, personalized, and capable of supporting monitoring, simulation, and decision-making?
\end{quote}

\section{Research Motivation}
The motivation of this thesis comes from the need to move from fragmented eHealth applications toward integrated and intelligent healthcare ecosystems. A Digital Twin-based approach can provide a structured way to combine real-world data, virtual representation, simulation, artificial intelligence, and decision services. This is particularly relevant in healthcare because clinical decisions depend not only on current measurements, but also on patient history, context, uncertainty, and possible future trajectories.

\section{Objectives}
The main objective of this thesis is to study the state of the art of Digital Twin technology in eHealth and to propose a general conceptual architecture for eHealth Digital Twin systems. The specific objectives are:

\begin{itemize}[leftmargin=*]
    \item to review the foundations of eHealth and digital health systems;
    \item to study the evolution, definitions, components, and maturity levels of Digital Twins;
    \item to analyze the use of Digital Twins in healthcare from a general eHealth perspective;
    \item to study major healthcare application domains, including personalized medicine, remote monitoring, hospital processes, public health, and cardiovascular applications;
    \item to compare existing architectures, technologies, applications, and limitations;
    \item to identify major challenges related to interoperability, security, privacy, ethics, and validation;
    \item to propose a conceptual multi-layer architecture for Digital Twin-based eHealth systems.
\end{itemize}

\section{Contributions}
The main contributions of this thesis are:

\begin{itemize}[leftmargin=*]
    \item a structured synthesis of Digital Twin concepts and their relevance to eHealth;
    \item a classification of healthcare Digital Twin applications from a general perspective;
    \item an extended discussion of cardiovascular Digital Twins as an advanced clinical application area;
    \item a comparative study of selected research papers and architectures;
    \item a proposed conceptual architecture for a Digital Twin-based eHealth system;
    \item a discussion of implementation challenges and future research directions.
\end{itemize}

\section{Thesis Outline}
This thesis is organized into five main parts. Part I introduces the research context, problem statement, objectives, and contributions. Part II presents the background on eHealth and Digital Twin fundamentals. Part III provides the literature review methodology and comparative state of the art. Part IV proposes a conceptual architecture for Digital Twin-based eHealth systems. Finally, Part V concludes the thesis and presents future work.

\clearpage
